\section{Introduction}
\label{sec:intro}

Data has become a central production factor of the modern economy, and trading is a primary means by which its value is unlocked and redistributed~\citep{analytics2016age,richter2019data}. 
A growing number of platforms---such as Snowflake Data Marketplace~\citep{SnowflakeDocs_MarketplaceAbout}, AWS Data Exchange~\citep{AWSDocs_DataExchangeLanding}, and Dawex~\citep{Dawex_Web_DataMarketplaceModel}---now intermediate the exchange of data products. 
Yet, in practice, data markets remain far less active than their potential would suggest~\citep{Koutroumpis2017UnfulfilledPotential,OECD2021ValueOfData,EuropeanParliament2023DataActPressRelease}, and a widely cited cause is the absence of incentive structures that reward all parties whose data ultimately creates value~\citep{OECD_EASD_2019,azcoitia2022survey,bauer2024designing}.

A substantial body of work studies incentive designs in data markets through auctions~\citep{ghosh2011selling,zhang2020selling,agarwal2024towards} and pricing~\citep{chen2019towards,mi2025multi}.
%, and federated learning~\citep{li2023martfl,yu2020fairness}. 
%However, ... (what is the problem)
However, most works treat data as physical goods, overlooking the \emph{replicable} nature of data:
% A defining feature that sets data apart from physical goods, however, is its \emph{replicability}:
once acquired, a buyer can process the data into a new product and resell it to multiple downstream buyers, so a single source can give rise to a cascade of resales~\citep{JonesTonetti2020NonrivalryData,shiller2013digital,biswas2021incentive}, naturally forming a tree structure of data flow. 
Because an upstream seller captures no value from these downstream resales, its incentive to trade in the first place is weakened.
%(因此如何建议seller在resales时的incentive是解决数据交易的关键问题。important.)
% Replicability also determines the shape these transactions take: a data product can serve many buyers simultaneously---a public machine learning dataset can feed countless model trainers---whereas a buyer rarely purchases the same data twice. 
% Data transactions thus naturally organize into a \emph{tree} rooted at the original source.
Recent work proposes \emph{profit reallocation}---redistributing part of downstream revenue back to upstream contributors---as a remedy for this incentive problem~\citep{xin2025tbds,ma2026incentivizing}.
Yet these analyses, like sequential-trade models more broadly~\citep{manea2018intermediation,condorelli2017bilateral}, assume that each seller resells to at most one buyer.
The design and effects of profit reallocation on the tree-structured markets that replicable data naturally induces remain unexplored.

To close this gap, we develop the first principled framework for profit reallocation on \emph{tree-structured} data markets.
We model data trading as a sequential game on a rooted tree in which the root player originates the data and every player may process and resell its data product to multiple children (\cref{sec:model}). 
On top of this game we formalize a \emph{profit reallocation mechanism} (PRM): whenever a trade occurs, part of its payment is redistributed to the buyer's ancestors by a budget-feasible redistribution rule---the total amount redistributed never exceeds the payment---so every player benefits from the trade.

Given a PRM, the first challenge is to understand how strategic players will behave, \ie, to compute the induced equilibrium; throughout we adopt the notion of \emph{perfect Bayesian equilibrium}, which demands optimality at every information set.
Branching makes this computation formidable: a seller prices data product to \emph{every} child, so its pricing action is a vector whose dimension equals its number of children, and its utility couples the entire subtree beneath it---a naive evaluation is exponential in the tree size.
Fortunately, a key structural insight dissolves this difficulty: under the Markovian valuation process, sibling subtrees are conditionally independent given the seller's valuation, so the seller's pricing problem \emph{separates} across its children.
This subtree decomposition collapses equilibrium strategies to single-variable functions and lets a seller optimize each child's price separately, eliminating any joint optimization over the price vector.
Building on it, we design a polynomial-time exact algorithm for discrete valuations and a fully polynomial-time approximation scheme (FPTAS) for continuous valuations (\cref{sec:algo}), turning the evaluation of any PRM into a tractable computation.

Beyond computation, we examine whether data trading actually benefits from profit reallocation, and answer in the affirmative.
We prove that, under mild assumptions, any budget-feasible PRM weakly expands the set of valuation profiles under which each trade occurs, relative to the baseline mechanism that does not reallocate profits; and any budget-balanced PRM, for which payments are fully redistributed rather than partly withheld by the mechanism, additionally weakly improves social welfare (\cref{sec:theoretic}).
Here, too, branching is the crux of the difficulty: rewarding a seller more from \emph{one} child's subtree could in principle distort its prices to the \emph{other} children, breaking trades elsewhere in the tree. We rule this out by carefully interpolating between PRMs so that adjacent PRMs differ on only one subtree at a time.
Finally, our experiments (\cref{sec:exp}) confirm that the advantage of profit reallocation is substantial, persists even when the theoretical assumptions are violated, and grows with the depth and branching of the tree.

In summary, our contributions are:
\begin{itemize}[left=0em]
\item A tree-based data trading game that captures one-to-many resale, together with a general, budget-feasible profit reallocation mechanism on it (\cref{sec:model});
\item Efficient equilibrium computation---a polynomial-time exact algorithm and an FPTAS---that makes any PRM evaluable (\cref{sec:algo});
\item A theoretical guarantee that any budget-feasible PRM weakly expands trade, and any budget-balanced PRM additionally weakly improves welfare, over the baseline mechanism on trees (\cref{sec:theoretic});
\item Experiments confirming that the advantage is robust to violations of the theoretical assumptions and strengthens as the tree size grows (\cref{sec:exp}).
\end{itemize}
Together, these results suggest that deploying profit reallocation can meaningfully invigorate real-world data markets.\footnote{Owing to space constraints, a broader discussion of related work is deferred to \cref{sec:related}.}

% Data has become a central production factor of the modern economy, and trading is a primary channel through which its value is unlocked and redistributed~\citep{analytics2016age,richter2019data}. A growing number of platforms---such as Snowflake Data Marketplace~\citep{SnowflakeDocs_MarketplaceAbout}, AWS Data Exchange~\citep{AWSDocs_DataExchangeLanding}, and Dawex~\citep{Dawex_Web_DataMarketplaceModel}---now intermediate the exchange of data products. Yet, in practice, data markets remain far less active than their potential would suggest~\citep{Koutroumpis2017UnfulfilledPotential,OECD2021ValueOfData,EuropeanParliament2023DataActPressRelease}, and a widely cited cause is the absence of incentive structures that reward all parties whose data ultimately creates value~\citep{OECD_EASD_2019,azcoitia2022survey,bauer2024designing}.

% A substantial body of work studies incentives in data markets through auctions~\citep{ghosh2011selling,zhang2020selling,agarwal2024towards}, pricing~\citep{chen2019towards,mi2025multi}, and federated-learning~\citep{li2023martfl,yu2020fairness}. A defining feature that sets data apart from physical goods, however, is its \emph{replicability}: once acquired, a buyer can process the data into a new product and resell it to multiple downstream buyers, so a single source can give rise to a cascade of resales~\citep{JonesTonetti2020NonrivalryData,shiller2013digital,biswas2021incentive}. Because an upstream seller captures no value from these downstream resales, its incentive to trade in the first place is weakened. Two recent works are most relevant here. \citet{xin2025tbds} introduce a transaction-based mechanism that pays upstream contributors a fixed share of downstream revenue, but their treatment is largely architectural and lacks an equilibrium analysis of the induced strategic behavior. \citet{ma2026incentivizing} provide the first systematic game-theoretic study of such \emph{profit reallocation mechanism}, but they assume that each seller can resell to at most one buyer thus the trades only form a chain. This assumption is restrictive: as digital goods, data can be duplicated and used by many buyers by nature. For example, a public machine learning dataset can feed uncountable model trainers. Such one-to-many resale naturally induces a \emph{tree} structure of transactions that the chain model cannot capture.

% To close this gap, we study profit reallocation on \emph{tree-structured} data markets. We model data trading as a sequential game on a rooted tree in which the root originates the data and every player may process and resell its data product to multiple children (\cref{sec:model}). On top of this game we formalize a \emph{profit reallocation mechanism} (PRM): a budget-feasible rule that, whenever a trade occurs, redistributes part of its payment to the buyer's ancestors. Our formulation only requires budget feasibility and strictly generalizes the fixed-ratio rules of prior work~\citep{xin2025tbds}, yielding a rich and economically meaningful design space.

% Given a PRM, a natural question is how strategic players will behave, \ie, how to compute the induced equilibrium; throughout we use the strong notion of \emph{perfect Bayesian equilibrium}, which demands optimality at every information set.
% The tree is qualitatively harder than the chain case considered by \citet{ma2026incentivizing}: on a chain each seller sets a single price, whereas on a tree a seller prices \emph{every} child, so its pricing action is a vector whose dimension is its number of children, and its utility couples the entire subtree beneath it---a naive evaluation is exponential in the tree size.
% Our key structural result is that, under the Markovian valuation process, sibling subtrees are conditionally independent given the seller's valuation, so the seller's pricing problem \emph{separates} across its children. This subtree decomposition collapses equilibrium strategies to single-variable functions and, crucially, lets a seller optimize each child's price separately, with no joint optimization over the price vector. Building on it, we design a polynomial-time exact algorithm for discrete valuations and a fully polynomial-time approximation scheme (FPTAS) for continuous valuations (\cref{sec:algo}), turning the evaluation of any PRM into a tractable computation.

% We further ask whether data trading actually benefits from profit reallocation. We prove that, under mild assumptions, any budget-feasible PRM weakly expands the set of valuation profiles under which each trade occurs, relative to the baseline mechanism that does not reallocate profits; and any budget-balanced PRM, for which the payments are fully redistributed rather than partly withheld by the mechanism, additionally weakly improves social welfare (\cref{sec:theoretic}). Here the tree case departs decisively from the chain case in \citet{ma2026incentivizing}: changing how a seller is rewarded from \emph{one} child subtree must be shown not to perturb its pricing to its \emph{siblings}---an obstacle absent on the chain, which we resolve by carefully interpolating PRMs such that adjacent PRMs differ on only one subtree at a time. Finally, our experiments (\cref{sec:exp}) show that the advantage of profit reallocation is substantial and persists even when the theoretical assumptions are violated, and that it grows with the depth and branching of the tree.

% In summary, our contributions are:
% \begin{itemize}[left=0em]
% \item A tree-based data trading game that captures one-to-many resale, together with a generalized profit reallocation mechanism on it (\cref{sec:model});
% \item Efficient equilibrium computations---a polynomial-time exact algorithm and an FPTAS---that makes any PRM evaluable (\cref{sec:algo});
% \item A theoretical guarantee that any budget-feasible PRM weakly expands trade, and any budget-balanced PRM additionally weakly improves welfare, over the baseline mechanism on trees (\cref{sec:theoretic});
% \item Experiments confirming that the advantage is robust to violations of the theoretical assumptions and strengthens as the tree size grows (\cref{sec:exp}).
% \end{itemize}
% Together, these results suggest that deploying profit reallocation can meaningfully invigorate real-world data markets.\footnote{Owing to space constraints, a broader discussion of related work is deferred to \cref{sec:related}.}